\section{Memory}
\label{sec:memory}

\subsection{A memory of verified answers}
\label{sec:mem:verified}
The context manager keeps conversations in SQLite and assembles each prompt
within a context cap, with cache-affinity so that the engine's prefix store
is hit. A retrieval memory holds earlier answers and offers them to later
questions, under one rule: \emph{only verified answers are stored}. A cache
that keeps every successful response would, by
Proposition~\ref{prop:distill}, feed unverified text back as if it were
knowledge. A hit is dropped when its stored question differs from the query
by at most two content words with a negation cue or an antonym pair on one
side only, so ``is caching enabled'' does not retrieve the stored answer to
``is caching disabled'' (a guard after vLLM Semantic Router
\cite{semanticrouter2026}). Retrieved memory enters the prompt with its
provenance and, like any retrieved text, as data.

\subsection{A memory of earlier conversations}
\label{sec:mem:recall}
The second memory holds what was said, so that a fact told in one conversation
can be found in a later one. It makes no model call.

\emph{Writing.} Each finished turn of a stored conversation becomes a note
(the question and the answer, reasoning removed, the answer clipped to
1{,}200 characters), and every third turn also a window over the last five.
A greeting, a turn of under 30 characters and an answer that declined are not
stored. A message that opens with ``Remember that'' is kept as a note that
the size limit never removes.

\emph{Reading.} Notes are read only when the question speaks in the first
person or points back (``my'', ``last time'', ``you said''), and never when
passages from a knowledge base are already in the prompt: on a CPU every
thousand characters of prompt is seconds of reading, so an impersonal question
pays nothing. Candidates come from BM25 over the notes and, when the
embedding server runs, from the cosine of note vectors, fused by reciprocal
rank. A note that answers the opposite question is dropped by the polarity
guard of Section~\ref{sec:mem:verified}. The tail below the largest drop in
relevance is cut. What remains is ordered by relevance, kind and age (a
half-life of 30 days), cleared of duplicates, and cut to three notes and
1{,}500 characters. Two notes are duplicates when their word sets overlap by
0.9 \emph{and} their numbers are bound to the same words, so ``budget 10K''
and ``budget 12K'' stay apart. The notes enter the prompt delimited as data.

\emph{Whose memory.} A private request, an experiment and a chat room neither
write nor read this memory; what is said in a room is shared by its members.
Deleting a conversation deletes its notes. A person can read every note,
delete one or all, add one, and merge repeated ones. At most 5{,}000 notes
are kept; those never recalled go first, oldest first.

The extraction, the gate, the elbow and the number-aware duplicate test follow
the conversation memory of vLLM Semantic Router \cite{semanticrouter2026}.
The polarity guard, the exclusion of rooms and private chats, a size limit
that uses the access counts, and the person's view of every note are added
here; a model-written rewrite of the query is left out. A live check is in
Section~\ref{sec:eval:recall}.

\subsection{The audit ledger as self-knowledge}
\label{sec:mem:ledger}
The ledger is the system's memory of its own performance: gate decisions,
verifier verdicts and audit outcomes per domain. It is what the calibrated
gate is fitted from (Section~\ref{sec:gate:labels}), and capabilities with no
audited data are reported as unmeasured, never as zero
(Proposition~\ref{prop:audit}). Verified deliberate traces are also the raw
material of distillation into the quick model, under the promotion rule of
Proposition~\ref{prop:distill}; training needs a GPU and was not run.

\subsection{Memory tokens}
\label{sec:mem:tokens}
Memorilla \cite{memorilla2026} compresses a document collection into a few
vectors that a frozen decoder reads in place of the embeddings of placeholder
tokens. Our port is a 31.5\,M-parameter module (sixteen slots, two
cross-attention blocks, eight heads, a projection from the embedder's width
to the decoder's) written in Candle, with forward-pass parity to the
unmodified PyTorch reference within $1.9\times10^{-6}$ on three
configurations. The engine admits requests whose placeholder positions carry
supplied vectors, and a trainer fits the module to a knowledge base through a
differentiable Qwen3 decoder (logits within $2\times10^{-5}$ of a reference
implementation). Document keys and values are prepared once per collection. The
\texttt{memory} chat mode uses the module alone and \texttt{hybrid} adds it to
retrieved passages; answers read from memory tokens are decoded greedily, as
the reference evaluates its own.

The module needs pre-training that a CPU cannot do, and the one module
fine-tuned here \emph{collapsed}: all sixteen memory vectors are one vector,
whatever the question (Section~\ref{sec:eval:memory}). Two measures follow.
The trainer now uses the reference's from-scratch recipe (learning rate
$10^{-4}$, sixteen examples per step, five epochs, a warm-up from zero over a
tenth of the steps, dropout 0.1 at the reference's four sites, a token-level
loss). And the trainer computes at every validation the mean cosine between
the slots of one question and between the memories of different questions,
saves both beside the weights, and orders checkpoints by whether the memory
depends on the question, then by the gain in claim accuracy over no memory,
and only then by validation loss. A module whose
memories for different questions have a cosine of 0.98 or more is not used:
the automatic mode answers from passages, and an explicit request for memory
falls back to passages and says why. A module saved without these statistics
is probed once with three fixed questions. The modes stay off by default.
